\documentclass[11pt]{article}
\usepackage[margin=0.85in]{geometry}
\usepackage{amsmath,amssymb,bm,booktabs,array,hyperref,xcolor,listings}
\hypersetup{colorlinks=true,linkcolor=blue,urlcolor=blue,pdftitle={Frozen-self-energy charge response and LAPW stress},pdfauthor={eDMFT implementation notes}}
\setlength{\emergencystretch}{2em}
\lstset{basicstyle=\ttfamily\small,breaklines=true,columns=fullflexible,frame=single}
\newcommand{\Tr}{\operatorname{Tr}}
\newcommand{\Herm}{\operatorname{Herm}}
\newcommand{\Sym}{\operatorname{Sym}}
\newcommand{\dd}{\mathrm d}
\newcommand{\ii}{\mathrm i}
\title{Frozen-self-energy charge response\\and the projector contribution to LAPW stress}
\author{Implementation and validation note for DFT+eDMFT}
\date{8 November 2026}
\begin{document}
\maketitle

\noindent\textbf{Purpose.} Obtain the strain derivative of the Kohn--Sham
potential without two converged strained calculations and without an impurity
solve. The immediate application is the chain-rule contribution from the
potential-dependent, normalized, finite-band-window projector. This note
distinguishes a tested complete charge-map derivative from the less accurate
stored-band approximation. The correction remains opt-in.

\section{What is differentiated, and what is kept fixed?}
Let $\eta$ be hydrostatic linear strain, with deformation gradient
$F_\eta=(1+\eta)I$, determinant $J_\eta=(1+\eta)^3$, and reference cell
volume $\Omega$. Fractional nuclear positions, muffin-tin radii, local
orbital axes, impurity self-energy on its supplied Matsubara mesh, and the
double-counting contribution contained in that input are held fixed in
the present tests. Band-window endpoints are fixed integer indices.

Use material coordinates $\mathbf x$ in the central cell:
\begin{equation}
 q_\eta(\mathbf x)=J_\eta n_\eta(F_\eta\mathbf x),\qquad
 v_\eta(\mathbf x)=V_{{\rm KS},\eta}(F_\eta\mathbf x).
 \label{eq:pullback}
\end{equation}
Here $n$ is the physical total electron density, including core charge;
$V_{\rm KS}$ is the local Kohn--Sham potential. A subscript 1 denotes
$\partial_\eta|_0$. Equation~\eqref{eq:pullback} gives
\begin{equation}
 q_1=n_1+3n_0+\mathbf x\cdot\nabla n_0,\qquad
 v_1=V_1+\mathbf x\cdot\nabla V_0.
 \label{eq:advection}
\end{equation}
The derivative $V_1$ on the right is at fixed physical position; $v_1$
is at fixed material position. Confusing the two loses the advection term.
For general infinitesimal strain $E$, replace $3$ by $\Tr E$ and
$\mathbf x\cdot\nabla$ by $(E\mathbf x)\cdot\nabla$.

The first implementation constrains total particle number. A Mott-insulator
calculation at fixed chemical potential is a different response problem;
the complete WIEN driver does not silently substitute fixed $N$ for that
ensemble. The analytic band-space action exposes both choices.

\section{The linear screening equation}
Write the linearized density-to-potential map in material coordinates as
\begin{equation}
 v_1=v_{\rm aff}+Kq_1.
 \label{eq:potential}
\end{equation}
$v_{\rm aff}$ is the potential derivative when $q_\eta=q_0$: nuclei and
cell are strained, while physical density is transported affinely.
$K$ is the central Hartree plus exchange-correlation kernel. In the numerical
implementation LAPW0 evaluates its action, including nonspherical,
interstitial and GGA contributions; no explicit dense kernel is built.

The frozen-self-energy electronic response satisfies
\begin{equation}
 q_1=b+\chi_{\rm fr}v_1,
 \qquad
 s=b+\chi_{\rm fr}v_{\rm aff},
 \qquad
 \boxed{(I-\chi_{\rm fr}K)q_1=s.}
 \label{eq:screen}
\end{equation}
$b$ is the explicit geometrical response of the one-electron charge map,
including basis and reconstruction changes, at fixed material potential.
The source $s$ is obtained together, so its strongly cancelling pieces need
not be evaluated separately. One screened approximation is $q_1\simeq s$,
followed by $v_1\simeq v_{\rm aff}+Ks$. It is not sufficient in these tests.

\subsection*{Why this is not a full interacting DMFT susceptibility}
In a fixed central orthonormal representation,
\begin{align}
 G_k(z)&=[z+\mu-H_k-U_k\widetilde\Sigma(z)U_k^\dagger]^{-1},\\
 \delta G_k&=G_k\left[\delta H_k-\delta\mu I
  +\delta U_k\widetilde\Sigma U_k^\dagger
  +U_k\widetilde\Sigma\delta U_k^\dagger\right]G_k.
 \label{eq:dg}
\end{align}
$z=\ii\omega_n$, $\widetilde\Sigma$ is the local self-energy after the
implemented double counting, and $U_k$ embeds local correlated orbitals.
There is no $U_k\delta\widetilde\Sigma_{\rm imp}U_k^\dagger$ in this
frozen-$\Sigma$ calculation. The physical DMFT susceptibility would require
that term and the impurity vertex. The screening in Eq.~\eqref{eq:screen}
uses the DFT kernel $K$, not an assumed Hubbard-$U$ RPA kernel.

Computing $\delta G$ here supplies $q_1$ and hence $v_1$. It does not mean
adding an independent $\delta G$ term to a stationary functional derivative.
Stationarity at fixed projectors does not remove the explicit chain rule
when the implemented projector itself depends on $V_{\rm KS}$.

\section{Finite window and normalization must both respond}
Let $\phi$ collect the untruncated localized trial orbitals, and let
$\Pi_{W,k}$ project onto the chosen KS bands at each $k$. Then
\begin{equation}
 U_k=\Pi_{W,k}\phi_k X,\qquad
 S=\Sym\sum_k w_k\phi_k^\dagger\Pi_{W,k}\phi_k,\qquad X=S^{-1/2}.
 \label{eq:projector}
\end{equation}
The operation $\Sym$ is the actual orbital-symmetry and equivalent-site
average used by the code. The $k$ weights $w_k$ are constant and sum to one.
For several impurity blocks these expressions apply with the implemented
block and site bookkeeping. Inactive orbital channels retain the code's
identity normalization; they must not acquire an artificial derivative.

For a perturbing potential matrix $D_k$ in the central KS basis, and
indices $n\in W$, $a\notin W$,
\begin{equation}
 (\delta\Pi_{W,k})_{an}=\frac{(D_k)_{an}}{\epsilon_{kn}-\epsilon_{ka}},
 \qquad
 (\delta\Pi_{W,k})_{na}=(\delta\Pi_{W,k})_{an}^*.
 \label{eq:pi}
\end{equation}
Within-$W$ and within-complement blocks vanish for this spectral-projector
derivative. A window boundary through a degeneracy is rejected, not
regularized. Smooth fixed-index selection still requires an isolated
window subspace.

In the frozen central LAPW basis,
\begin{align}
 \delta S&=\Sym\sum_k w_k\phi_k^\dagger\delta\Pi_{W,k}\phi_k,\\
 \delta U_k&=\delta\Pi_{W,k}\phi_kX+\Pi_{W,k}\phi_k\delta X.
 \label{eq:du}
\end{align}
If $S=V\operatorname{diag}(s_i)V^\dagger$, the exact Fr\'echet derivative is
\begin{equation}
 (V^\dagger\delta X V)_{ij}
 =-\frac{(V^\dagger\delta S V)_{ij}}
 {\sqrt{s_i}\sqrt{s_j}(\sqrt{s_i}+\sqrt{s_j})}.
 \label{eq:dx}
\end{equation}
No projector pseudoinverse is needed: the exported $X$ is authoritative.
The saved normalization can differ slightly from the newly summed $S$;
directional endpoint tests anchor to the saved central $X$ consistently.

The stored-band implementation uses only available complement bands in
Eq.~\eqref{eq:pi}. It also holds central radial augmentation fixed.
The complete charge-map derivative below includes the response missing
from both restrictions, while keeping the same physical DMFT window.

\section{Matsubara summation and particle-number constraint}
For one spin component define
$H_{\infty,k}=H_k+U_k\widetilde\Sigma_\infty U_k^\dagger-\mu I$,
$G_{\infty,k}(z)=(z-H_{\infty,k})^{-1}$, and
$f_T(x)=(e^{x/T}+1)^{-1}$. The implemented density prescription is
\begin{equation}
 \rho_k=g_s w_k\left[f_T(H_{\infty,k})+
 2T\sum_{n\in{\cal M}} a_n\Herm\{G_k(\ii\omega_n)-G_{\infty,k}(\ii\omega_n)\}\right].
 \label{eq:rho}
\end{equation}
$g_s=2$ in the tested nonmagnetic calculation. $\mathcal M$ is the kept
positive Matsubara mesh and $a_n$ are the actual production quadrature
weights. The native export supplies them, $T$, $\mu$, and the effective
self-energy matrices in Ry. The action differentiates Eq.~\eqref{eq:rho},
including the static subtraction and exact Fermi-matrix derivative. It
does not square a quadrature-weighted Green function: $a_n$ appears once
after the two resolvents in Eq.~\eqref{eq:dg} are multiplied.

In an eigenbasis of $H_\infty$, the static derivative is
\begin{equation}
 [L_f(H_\infty;\delta H_\infty)]_{ij}
 =f_T[\varepsilon_i,\varepsilon_j](\delta H_\infty)_{ij},\quad
 f_T[x,y]=\frac{f_T(x)-f_T(y)}{x-y},\quad f_T[x,x]=f_T'(x).
\end{equation}
Let $\delta\rho^{(0)}$ be the fixed-$\mu$ response and $\rho^{(\mu)}$
the response to a unit increase in $\mu$. Fixed total particle number gives
\begin{equation}
 \delta\mu=-\frac{\sum_k\Tr\delta\rho_k^{(0)}}
 {\sum_k\Tr\rho_k^{(\mu)}},\qquad
 \delta\rho_k=\delta\rho_k^{(0)}+\delta\mu\,\rho_k^{(\mu)}.
\end{equation}
A nearly zero denominator requires an explicit ensemble decision. The
analytic action refuses to divide by it. A uniform potential shift at
fixed $N$ gives $\delta U=0$, $\delta\mu=\delta V$, and $\delta\rho=0$.

\section{Complete directional derivative without an impurity solve}
Define ${\cal F}_\eta[q]$ as \emph{one} frozen-$\Sigma$ charge evaluation,
pulled back to the reference cell. Each evaluation consists of LAPW0,
LAPW1, dmft2 with refreshed normalization and chemical potential, lcore,
and the standard interstitial charge normalization. No mixer history or
impurity solve is used. Then
\begin{align}
 s&\simeq\frac{{\cal F}_{+h}[q_0]-{\cal F}_{-h}[q_0]}{h_+-h_-},\\
 A x&\equiv\chi_{\rm fr}Kx
 \simeq\frac{{\cal F}_0[q_0+t x]-{\cal F}_0[q_0]}{t}.
 \label{eq:jvp}
\end{align}
Actual rounded lattice strains $h_\pm$ are used. The central baseline is
recomputed with the same program path and normalization as each action.
This avoids dividing a small central mismatch by $t$.

The second line is a numerical action of a linear-response operator at the
central cell, not a converged strained calculation. It includes radial
response and the complete LAPW eigenvector response implicitly: computing
the occupied/window eigenvectors at the perturbed potential lets all basis
states participate without writing every unoccupied band.

Plain mixing can diverge. With residual $r_j=s+Aq_j-q_j$, the implementation
uses Anderson history:
\begin{equation}
 \gamma=\arg\min_\gamma\|r_j-\Delta R\gamma\|_2,\qquad
 q_{j+1}=q_j+\beta r_j-(\Delta Q+\beta\Delta R)\gamma.
\end{equation}
Here $\Delta Q$ and $\Delta R$ contain differences of recent iterates and
residuals. Only this small history problem uses SVD. The returned iterate
must have a measured residual below tolerance; otherwise the driver exits
with a nonconvergence status. Tolerances refer to the stored density-vector
norm, not to a universal physical norm; stress and step convergence are
checked independently.

\subsection*{LAPW-specific safeguards}
\begin{itemize}
\item The source and response runs use exactly the central integer G set,
through the existing serial WIEN fixed-G reference markers. Fixed band
count alone does not ensure a differentiable basis.
\item Real versus complex LAPW1 is selected from its unit-5 definition,
not from the mere presence of an \texttt{in1c} file. Empty optional files
can exist in real-vector cases.
\item The spherical \texttt{clmval} channel contains
$\sqrt{4\pi}\,r^2 n$, whereas \texttt{clmsum}/\texttt{clmcor} contain
$4\pi r^2 n$. Only the $l=m=0$ valence channel needs the extra
$\sqrt{4\pi}$ factor before adding core density. Higher harmonics are
unchanged.
\item For radial density coefficients $c(r)=r^2n_{lm}(r)$, affine transport
is $c_\lambda(r)=c_0(r/\lambda)/\lambda$; pullback is
$\lambda c_\lambda(\lambda r)$. Fourier density star sums acquire
$\lambda^{-3}$ and $\lambda^3$, respectively. Cubic interpolation on the
logarithmic radial mesh is used. Boundary extrapolation is small but must
be checked by halving $h$; density at $R_{\rm MT}$ is not assumed zero.
\item The analytic signed density matrix is reconstructed through the
existing dmft2 radial/angular/FFT machinery. All signed eigenmodes are
retained. The input includes eigenvectors to reject a mismatched gauge,
including rotations within degenerate subspaces. Energies printed during
this reconstruction are not response energies.
\end{itemize}

\section{Projector stress contraction and its limits}
Let $R_k^{W}$ and $R_k^{X}$ be the already implemented adjoint matrices for
the window and normalization chain rules, defined operationally by
\begin{equation}
 \delta\Gamma_{\rm proj}^{V}
 =\Re\sum_k\Tr[(R_k^{W}+R_k^{X})\,\delta v_k^{\rm op}].
\end{equation}
Here $\delta v_k^{\rm op}$ is the matrix of the \emph{material} local-potential
derivative in the central KS representation. The exported adjoints already
contain the code's spin, frequency, and $k$ weights; none is applied again.
The correction to mean tensile stress is
\begin{equation}
 \boxed{\Delta\bar\sigma=
 \frac{1}{3\Omega}\Re\sum_k\Tr[(R_k^{W}+R_k^{X})(v_{\rm aff}+Kq_1)_k^{\rm op}].}
\end{equation}
Multiply Ry/Bohr$^3$ by $14710.507641078$ for the GPa conversion used in
these comparison archives. Compressive pressure has the opposite sign.
The tabulated values are already hydrostatic means; do not divide by three
again. For tetragonal FeSe this is the trace divided by three, not any
individual diagonal component.

The diagnostic contraction uses a frozen-central-basis LAPW1 matrix probe
and a complete central dual basis, with separate MT advection and
interstitial material-coordinate treatment. This probe is a reference
integrator; it does not imply that the density response can freeze its
radial basis. The general production replacement would evaluate those
potential matrix elements directly.

The driver writes predicted physical endpoint potentials using
$n_\pm(\mathbf r)=J_\pm^{-1}[q_0\pm hq_1](F_\pm^{-1}\mathbf r)$ and two
LAPW0 calls. Subtracting the raw \texttt{vns} arrays alone is not the
material potential operator: sphere/interstitial and advection terms
must still be treated consistently.

\section{Validation results}
All results below are mean tensile stresses in GPa, using saved
frozen-$\Sigma$ finite-difference Klein references from the preceding audit.
Those reference totals were not recomputed by an impurity solver during this
work. Neither their densities nor potentials were inputs to the response
solver; they were used only for independent validation.

\begin{center}\small
\begin{tabular}{lrrrr}\toprule
Material & Previous analytic & Corrected & Numerical & Difference\\\midrule
SrVO$_3$ & +5.676884 & +2.690161 & +2.933836 & -0.243675 \\
FeSe & -0.116492 & -1.118667 & -1.049151 & -0.069516 \\
\bottomrule\end{tabular}\end{center}

\begin{center}\small
\begin{tabular}{lrrrr}\toprule
Material & Predicted correction & Relaxed reference & Difference & Actions\\\midrule
SrVO$_3$ & -2.98672306 & -2.98675885 & +3.58e-05 & 9 \\
FeSe & -1.00217579 & -1.00221547 & +3.97e-05 & 11 \\
\bottomrule\end{tabular}\end{center}

The final validation used a fixed integer G set, fixed band
counts (43 in SrVO$_3$, 66 in FeSe), and fixed windows $[12,31]$ and
$[21,44]$, respectively. The retained meshes have 75 and 140 k points.
The source strain was $h=2.5\times10^{-4}$ and the central density-probe
step $t=10^{-5}$. The screened-density residual norms were 0.01213 and
0.00893. Their differences from the relaxed-reference material-density
response were 0.01504 and 0.01352, versus reference norms 1062.99 and 997.37.

Halving both steps to $h=1.25\times10^{-4}$, $t=5\times10^{-6}$ gives:
\begin{center}\small
\begin{tabular}{lrr}\toprule
Material & Correction, half steps & Change\\\midrule
SrVO$_3$ & -2.98672795 & -4.88e-06 \\
FeSe & -1.00219128 & -1.55e-05 \\
\bottomrule\end{tabular}\end{center}

\subsection*{What exposed the incomplete approximation?}
Insert the \emph{known} relaxed-reference density derivative into the
right-hand side of Eq.~\eqref{eq:screen}; measure the residual against that
reference. Using the corrected fixed-G source, the density-vector norms are:
\begin{center}\small
\begin{tabular}{lrr}\toprule
Charge response used & SrVO$_3$ & FeSe\\\midrule
Stored bands, frozen radial basis & 170.264 (43 bands) & 426.406 (66 bands)\\
100 stored bands, frozen radial basis & 109.944 & ---\\
200 stored bands, frozen radial basis & 89.013 & ---\\
Complete LAPW response, frozen radial basis & 75.731 & ---\\
Complete LAPW and radial response & 0.0174 & 0.0562\\\bottomrule
\end{tabular}\end{center}
The complete frozen-radial test freezes the basis in \emph{both} LAPW1 and
dmft2 reconstruction. Freezing only LAPW1 would mix inconsistent radial and
coefficient representations. The results establish that adding a modest
number of complement bands alone does not fix the charge-response action.
Radial-basis response is significant in SrVO$_3$.

\subsection*{Checks passed}
\begin{itemize}
\item The exported-mesh Python density reproduces the production central
band density to $3.3\times10^{-15}$ in SrVO$_3$ and $9\times10^{-16}$ in FeSe.
\item Independent nonlinear spectral-projector endpoint differences verify
Eq.~\eqref{eq:dg} and the normalization derivative with quadratic step
convergence. A small standalone algebra fixture tests the same identity.
\item Fixed-$N$ conservation and uniform-potential gauge tests pass.
\item Signed density reconstruction retains negative modes; an intentionally
changed eigenvector gauge is rejected with a failing process status.
\item With optional hooks disabled, SrVO$_3$ charge output is identical to
the prior result, and FeSe differs by at most $1.1\times10^{-18}$ in the
stored density coefficients.
\item LAPW0 applied to the measured reference density derivative reproduces
the independently measured potential correction to approximately
$2\times10^{-5}$ GPa (SrVO$_3$) and $5\times10^{-5}$ GPa (FeSe).
\item Final fixed-G source and action runs preserve the G sets, band counts,
k ordering and isolated window boundaries. Early preliminary runs that
allowed G changes are superseded by these results.
\end{itemize}

The screening actions took about 169 seconds in SrVO$_3$ and 138 seconds in
FeSe on this laptop, with four MPI ranks for dmft2 and serial LAPW1.
Summed program times including source and central evaluations were about
218 and 172 seconds, before the separate stress matrix probe. These are
measured diagnostic costs, not a general scaling guarantee. There were no
impurity solves and no pairs of converged strained charge calculations.

\textbf{Remaining error.} The total-stress residuals of 0.244 and 0.070 GPa
are not explained by the numerical accuracy of this potential-response
calculation. That does not identify their source: residual analytic basis
approximations, the complement convergence of the stress adjoint itself,
and consistency of the functional finite-difference reference still need
separate checks. Agreement of the new correction with its reference must
not be reported as exact agreement of the complete stress.


\section{Files and practical use}
The reusable sources are in \texttt{src/python/}:
\begin{itemize}
\item \texttt{frozen\_sigma\_response.py}: analytic fixed-basis stored-band
action, including projector/window/normalization and Matsubara derivatives.
\item \texttt{frozen\_sigma\_screening.py}: matrix-free Anderson solver.
\item \texttt{wien\_response\_density.py}: density conventions and transport.
\item \path{wien_frozen_sigma_response.py}: complete WIEN charge-map
action, source construction, screening and potential output.
\item \path{frozen_sigma_projector_stress.py}: guarded hydrostatic contraction of the saved potential and projector-adjoint matrices.
\end{itemize}
Native hooks are in \texttt{src/dmft2/frozen\_response\_io.f90}, with optional
calls from \texttt{l2main.f90}. They are disabled unless requested. A
separate \texttt{dmft2\_response} executable is installed; the ordinary
installed dmft2 is not replaced by an automatic stress correction.

The complete driver takes a JSON configuration. Its workspace must be new.
It copies inputs and never modifies the reference calculation. A template is
included beside this note. Activate the usual environment, then run:
\begin{lstlisting}
source ~/activate
~/dbin/wien_frozen_sigma_response.py response_config.json
\end{lstlisting}
The configuration names the reference case, workspace, WIEN root and executable,
MPI launcher, and LAPW1/dmft2 definition files. Definitions must use relative
case paths; the driver rewrites the distributed-vector suffix for the serial
LAPW1 output. The initial complete driver supports nonmagnetic primitive
orthogonal cells with fixed $N$ and projector 5. It does not claim magnetic,
SOC, shear, or general-cell validation.

For the analytic action, first replay a central TOT dmft2 charge calculation
with recomputeEF=0 and the saved converged chemical potential:
\begin{lstlisting}
export EDMFT_RESPONSE_EXPORT=/absolute/work/central
mpirun -n 4 ~/dbin/dmft2_response dmft2.def
unset EDMFT_RESPONSE_EXPORT
~/dbin/frozen_sigma_response.py --export /absolute/work/central \
  --potential dV_band.npy --output /absolute/work/density_response
export EDMFT_RESPONSE_DENSITY_PREFIX=/absolute/work/density_response
mpirun -n 4 ~/dbin/dmft2_response dmft2.def
unset EDMFT_RESPONSE_DENSITY_PREFIX
\end{lstlisting}
Run the last command in an isolated copy: its \texttt{clmval} is a signed
response density and must not replace an ordinary converged density. The
matrix input has shape $(N_k,N_b,N_b)$ and units Ry per strain, in exactly
the exported central eigenvector gauge. $N_b$ includes all stored bands.

\section{What remains open}
The tested response can reproduce the potential-dependent projector
correction obtained from the saved relaxed endpoints. It does not by itself
prove that every remaining term in the stress functional is exact. The
residual total-stress errors remain explicit in the table. Their audit
should use the same ensemble, normalized projector, band/G sets, and
frozen-$\Sigma$ Klein comparison including its impurity-functional linear
correction. No values from MnO or magnetic MnTe are inferred from these
nonmagnetic metallic tests.

For speed, the next production step is an analytic radial response and a
complete complement action replacing the numerical charge-map tangent.
The stored-band implementation provides a useful building block and an
independent check, but the measured radial error rules out promoting it
unchanged. The present complete tangent is a validated, moderately priced
reference route for deriving that faster implementation.
\end{document}
